\documentclass[preprint,12pt]{elsarticle}

\usepackage[utf8]{inputenc}
\usepackage[T1]{fontenc}
\usepackage{hyperref}
\usepackage{graphicx}
\usepackage{booktabs}
\usepackage{amsmath}
\usepackage{amssymb}
\usepackage{float}
\usepackage{subcaption}

\hypersetup{
  colorlinks=true,
  linkcolor=blue,
  citecolor=blue,
  urlcolor=blue
}

\journal{SoftwareX}

\begin{document}

\begin{frontmatter}

\title{ChemIP: An open-source platform integrating MSDS, patent, and trade intelligence for chemical safety decision-making}

\author[inst1]{Yu Yong Kim\corref{cor1}}
\cortext[cor1]{Corresponding author}
\ead{ykim288@wisc.edu}

\affiliation[inst1]{organization={University of Wisconsin-Madison}, country={United States}}

\begin{abstract}
Chemical safety practitioners must routinely consult multiple disconnected government portals---each with its own interface, query syntax, and response format---to evaluate a single substance.
This fragmented workflow increases investigation time and introduces risks of information omission.
ChemIP is an open-source, self-hostable web platform that unifies nine Korean and international public data sources---covering MSDS safety data, patents, trade intelligence, pharmaceutical registries, chemical classification, and biomedical literature---into a single-query workflow with adapter-based API orchestration.
The platform is built with FastAPI and Next.js, supports local deployment for privacy-sensitive environments, and includes an optional LLM-assisted summary with deterministic rule-based fallback.
The source code, automated tests, and CI/CD pipeline are publicly available under the AGPL-3.0 license.
\end{abstract}

\begin{keyword}
chemical safety \sep MSDS \sep API orchestration \sep open-source \sep decision support \sep public data integration
\end{keyword}

\end{frontmatter}

% ============================================================
% Code Metadata (MANDATORY for SoftwareX)
% ============================================================

\begin{table}[H]
\caption{Code metadata.}
\label{tab:code_metadata}
\small
\begin{tabular}{@{}lp{9cm}@{}}
\toprule
\textbf{Nr.} & \textbf{Code metadata description} \\
\midrule
C1 & Current code version: v1.0.0 \\
C2 & Permanent link to code/repository: \url{https://github.com/yuyongkim/chemIP} \\
C3 & Permanent link to Reproducible Capsule: -- \\
C4 & Legal Code License: AGPL-3.0 \\
C5 & Code versioning system used: git \\
C6 & Software code languages, tools, and services used: Python~3.13, TypeScript~5, FastAPI, Next.js~16, React~19, SQLite, Ollama \\
C7 & Compilation requirements, operating environments \& dependencies: Python~3.11+, Node.js~20+, 34 Python packages, 9 npm packages; Windows/Linux/macOS \\
C8 & Link to developer documentation: \url{https://github.com/yuyongkim/chemIP/blob/main/README.md} \\
C9 & Support email for questions: ykim288@wisc.edu \\
\bottomrule
\end{tabular}
\end{table}

% ============================================================
\section{Motivation and Significance}
\label{sec:motivation}

Chemical safety management requires practitioners to cross-reference at least four categories of information when evaluating a substance: (1)~hazard classifications from Material Safety Data Sheets (MSDS); (2)~patent status for freedom-to-operate analysis; (3)~trade regulations and market dynamics; and (4)~pharmaceutical regulatory status when chemicals overlap with drug formulations.

In the Republic of Korea, these domains are managed by separate agencies---KOSHA (occupational safety), KIPRIS (patents), KOTRA (trade), MFDS (pharmaceuticals), KISCHEM (chemical safety), and NCIS/Korea Environment Corporation (substance classification)---each with independent web portals, query interfaces, and response formats.
A typical evaluation requires sequentially accessing each portal, re-entering search queries, and manually compiling results into a working document.
This fragmented workflow introduces three systemic risks: \emph{information omission} due to variable practitioner expertise, \emph{temporal overhead} from sequential portal access with redundant authentication, and \emph{cross-domain blindness} that prevents identifying connections between safety, patent, and market data.

The core problem is not a lack of public data, but \textbf{workflow fragmentation}: the information exists across well-maintained government services, yet no open tool integrates them into a unified lookup process.

Table~\ref{tab:comparison} compares ChemIP with existing platforms.
PubChem~\cite{pubchem_2023} provides comprehensive chemical properties but does not integrate MSDS operational data, patents, or trade intelligence.
SciFinder~\cite{scifinder} and Reaxys~\cite{reaxys} offer patent-integrated molecular search but are commercial products that do not cover Korean regulatory data.
OpenFDA~\cite{openfda} covers US drug safety without cross-domain integration.
ChemIP is, to the best of our knowledge, one of the few openly available tools that integrates MSDS, patent, trade, pharmaceutical, chemical classification, and literature data in a single self-hostable workflow.

\begin{table}[H]
\centering
\caption{Feature comparison of chemical information platforms.}
\label{tab:comparison}
\small
\begin{tabular}{@{}lccccccc@{}}
\toprule
\textbf{Feature} & \rotatebox{70}{\textbf{ChemIP}} & \rotatebox{70}{\textbf{PubChem}} & \rotatebox{70}{\textbf{SciFinder}} & \rotatebox{70}{\textbf{Reaxys}} & \rotatebox{70}{\textbf{OPS}} & \rotatebox{70}{\textbf{OpenFDA}} & \rotatebox{70}{\textbf{KOSHA Web}} \\
\midrule
Open-source          & \checkmark & \checkmark & -- & -- & \checkmark & \checkmark & -- \\
MSDS integration     & \checkmark & -- & partial & -- & -- & -- & \checkmark \\
Patent search        & \checkmark & -- & \checkmark & \checkmark & \checkmark & -- & -- \\
Trade intelligence   & \checkmark & -- & -- & -- & -- & -- & -- \\
Drug cross-ref       & \checkmark & partial & -- & -- & -- & \checkmark & -- \\
Substance classif.   & \checkmark & partial & -- & -- & -- & -- & -- \\
Multi-API orchestr.  & \checkmark & -- & -- & -- & -- & -- & -- \\
Self-hosted          & \checkmark & -- & -- & -- & -- & -- & -- \\
\bottomrule
\end{tabular}
\end{table}

% ============================================================
\section{Software Description}
\label{sec:software}

\subsection{Software Architecture}
\label{sec:architecture}

ChemIP follows a three-tier architecture (Fig.~\ref{fig:architecture}): a Next.js~16 frontend with React~19, a FastAPI backend with 30+ REST endpoints organized into seven route modules, and a data layer comprising SQLite databases and nine external APIs.

\begin{figure}[H]
\centering
\fbox{\parbox{0.92\textwidth}{\centering
\textbf{ChemIP System Architecture}\\[6pt]
\texttt{[Web Browser]} $\rightarrow$
\texttt{[Next.js 16 / React 19 --- port 7000]}\\[4pt]
$\downarrow$ \small{REST/JSON proxy}\\[4pt]
\texttt{[FastAPI Backend --- port 7010]}\\
\small{30+ endpoints $\cdot$ 5-layer middleware $\cdot$ 7 route modules}\\[4pt]
$\swarrow$ \hspace{0.8cm} $\downarrow$ \hspace{0.8cm} $\searrow$\\[4pt]
\parbox{2.8cm}{\centering\small\texttt{[External APIs]}\\KOSHA, KIPRIS,\\KOTRA (7), MFDS,\\OpenFDA, PubMed,\\Naver, KISCHEM,\\NCIS}
\hspace{0.2cm}
\parbox{2.8cm}{\centering\small\texttt{[SQLite DBs]}\\terminology.db\\(FTS5 index)\\patent\_index.db\\(optional)}
\hspace{0.2cm}
\parbox{2.8cm}{\centering\small\texttt{[Local LLM]}\\Ollama\\(optional)}
}}
\caption{Three-tier architecture. The frontend proxies API requests to the backend, which orchestrates nine external APIs, local databases, and an optional LLM. Components marked ``optional'' are not required for core functionality.}
\label{fig:architecture}
\end{figure}

\textbf{Adapter-based integration.}
Each external data source is wrapped in a dedicated adapter module (Table~\ref{tab:sources}) that handles authentication, request construction, response parsing (XML or JSON), and error normalization.
All adapters share a unified HTTP client with connection pooling (20 connections), exponential-backoff retry ($t_k = 0.5 \times 2^k$~s, max 3~retries), and automatic failover on transient errors (HTTP 502/503/504).
This design allows new data sources to be added without modifying existing adapters or route handlers.

The KOSHA MSDS adapter includes a fallback mechanism: when the primary API endpoint (\texttt{msds.kosha.or.kr}) is unreachable, requests are automatically rerouted to the \texttt{data.go.kr} proxy endpoint, ensuring continued access during service outages.

\begin{table}[H]
\centering
\caption{Integrated data sources and adapter characteristics.}
\label{tab:sources}
\small
\begin{tabular}{@{}llrl@{}}
\toprule
\textbf{Source} & \textbf{Format} & \textbf{LOC} & \textbf{Scope} \\
\midrule
KOSHA MSDS & XML & 98 & 16-section safety data (with fallback) \\
KIPRIS Patents & XML & 209 & Korean patent search/detail \\
KOTRA (7~APIs) & JSON/XML & 422 & Trade intelligence \\
MFDS & JSON & 70 & Korean drug approvals \\
OpenFDA & JSON & 30 & US drug labels \\
PubMed & JSON & 66 & Biomedical literature \\
Naver & JSON & 54 & Korean news search \\
KISCHEM & XML & 82 & Exposure symptoms, first-aid \\
NCIS & JSON & 78 & Substance classification \\
\bottomrule
\end{tabular}
\end{table}

\textbf{Graceful degradation.}
The system is designed to function when individual components are unavailable.
If the KOSHA API is down, cached MSDS data is served from the local database.
If the optional patent index database is not installed, patent search falls back to live KIPRIS API queries.
If the LLM server (Ollama) is offline, a deterministic rule-based report builder generates structured safety summaries from the same evidence bundle, with a confidence score $C = \min(0.95,\; 0.4 + 0.04 \times N_{\text{evidence}})$.

\textbf{Optional patent search backend.}
For deployments requiring large-scale patent analysis, ChemIP supports an optional pre-built patent index using Aho-Corasick multi-pattern matching~\cite{ahocorasick} over an FTS5-indexed SQLite database~\cite{sqlite_fts5}.
This component is not required for core platform functionality and is documented separately in the repository.

\textbf{Guide recommendation.}
A multi-criteria engine scores chemical--guide matches using weighted exact-match scoring: $S = 4|M_{\text{CAS}}| + 3|M_{\text{title}}| + 2|M_{\text{body}}|$, with deterministic 8-tuple ranking for reproducible ordering.

\subsection{Software Functionalities}
\label{sec:functionalities}

ChemIP provides six integrated workflows through a tabbed interface (Figs.~\ref{fig:main}--\ref{fig:detail}):

\begin{enumerate}
  \item \textbf{Substance safety lookup}: A single query (chemical name, CAS number, or English name) retrieves all 16 MSDS sections with GHS pictograms, cached locally after first retrieval.
  \item \textbf{Patent landscape}: Dual-track search combining real-time KIPRIS API queries with optional local index lookups. Results are classified as exclusion, usage, or mention by context analysis.
  \item \textbf{Trade intelligence}: Four analytical views (market news, entry strategy, price trends, fraud alerts) plus HS-code-based import restriction checking via KOTRA APIs.
  \item \textbf{Pharmaceutical cross-reference}: Aggregates Korean drug approvals (MFDS), US drug labels (OpenFDA), and biomedical literature (PubMed) for the queried substance.
  \item \textbf{Korean regulatory classification}: Aggregates KISCHEM exposure symptom and first-aid data with NCIS substance classification (toxic, restricted, accident-preparedness) and molecular data.
  \item \textbf{AI-assisted synthesis}: Collects evidence across all domains and generates a natural-language safety summary. When the LLM is unavailable, a rule-based fallback produces a structured report with source attribution.
\end{enumerate}

\begin{figure}[H]
\centering
\begin{subfigure}[t]{0.48\textwidth}
  \centering
  \includegraphics[width=\textwidth]{screenshots_en/cap_main.png}
  \caption{Home page with unified search.}
  \label{fig:main}
\end{subfigure}
\hfill
\begin{subfigure}[t]{0.48\textwidth}
  \centering
  \includegraphics[width=\textwidth]{screenshots_en/cap_chemical.png}
  \caption{Chemical detail with MSDS and GHS data.}
  \label{fig:detail}
\end{subfigure}
\\[8pt]
\begin{subfigure}[t]{0.48\textwidth}
  \centering
  \includegraphics[width=\textwidth]{screenshots_en/cap_patents.png}
  \caption{Patent search interface.}
  \label{fig:patents}
\end{subfigure}
\hfill
\begin{subfigure}[t]{0.48\textwidth}
  \centering
  \includegraphics[width=\textwidth]{screenshots_en/cap_trade.png}
  \caption{KOTRA trade intelligence dashboard.}
  \label{fig:trade}
\end{subfigure}
\caption{ChemIP user interface. (a)~Unified search with autocomplete. (b)~Eight-tab chemical detail view. (c)~Patent search with optional local index backend. (d)~Trade analysis hub.}
\label{fig:screenshots}
\end{figure}

% ============================================================
\section{Illustrative Usage Scenarios}
\label{sec:examples}

The following scenarios illustrate how ChemIP consolidates multi-portal workflows into unified queries. These are representative usage examples, not formal evaluations.

\textbf{Scenario 1: New raw material safety review.}
A safety officer evaluates toluene (CAS:~108-88-3) as a candidate solvent.
A single search retrieves the full 16-section MSDS (GHS pictograms GHS02, GHS07, GHS08; TWA:~50\,ppm), three matched KOSHA safety guides, KIPRIS patent results, and NCIS classification data---information that would otherwise require visiting at least four separate portals.

\textbf{Scenario 2: Export regulatory assessment.}
A trade analyst checks methanol (CAS:~67-56-1) export requirements for Vietnam.
One query retrieves GHS classification, KOTRA entry strategy data, import restriction status, patent landscape, and trade fraud alerts, providing a consolidated regulatory view.

\textbf{Scenario 3: Pharmaceutical cross-reference.}
A researcher investigates salicylic acid (CAS:~69-72-7).
The drug workflow aggregates 23 MFDS-approved products, OpenFDA labels with relevant warnings, and recent PubMed articles, alongside the MSDS irritant classification---consolidating sources that span Korean and US regulatory databases.

% ============================================================
\section{Impact}
\label{sec:impact}

\subsection{Practical Utility}

ChemIP addresses a concrete workflow problem: the manual effort required to navigate between disconnected government data portals during substance evaluation.
By consolidating nine data sources into a single-query interface, the platform eliminates redundant search entry, provides local caching of previously retrieved data, and enables tab-based cross-domain navigation without portal switching.

The self-hosted architecture ensures that chemical inquiry patterns---which may reveal proprietary formulations or pending patent strategies---remain within the deployment environment.
This is a requirement in many industrial settings where substance evaluations involve trade-sensitive information.

\subsection{Software Quality and Reproducibility}

The platform includes 19 automated smoke tests covering all API endpoint categories, a GitHub Actions CI/CD pipeline that runs on every push to the main branch, and single-command startup scripts for both development and production environments.
The test suite uses monkeypatch-based mocking to ensure no external API calls during testing.

\subsection{Deployment Modes}

ChemIP supports three deployment configurations:
\begin{itemize}
  \item \textbf{Minimal}: Core search and MSDS lookup using the terminology database and live API calls. No additional setup required beyond API keys.
  \item \textbf{Standard}: Adds cached MSDS data, guide recommendations, drug cross-reference, and Korean regulatory lookups.
  \item \textbf{Full}: Adds the optional local patent index and LLM-assisted analysis.
\end{itemize}

\subsection{Limitations}

\begin{itemize}
  \item No formal user study has been conducted; usability evaluation with domain practitioners is planned as future work.
  \item Some workflows depend on external government APIs that may change specifications or experience outages, as observed with intermittent KOSHA service interruptions during development.
  \item The optional patent index requires substantial storage and is not necessary for core platform functionality.
  \item The LLM component (Qwen3:8b) is a general-purpose model; domain-specific fine-tuning could improve the relevance of generated analyses.
  \item Full deployment requires API keys from Korean government data portals (\texttt{data.go.kr}), though the platform operates in degraded mode without them.
\end{itemize}

% ============================================================
\section{Conclusions}
\label{sec:conclusions}

ChemIP provides an open-source, adapter-based integration platform that unifies nine public data sources---spanning MSDS safety data, patents, trade intelligence, pharmaceutical registries, chemical classification, and biomedical literature---into a single-query workflow for chemical safety practitioners.
The platform is designed for local deployment, supports graceful degradation when individual components are unavailable, and includes automated tests with CI/CD.
Future work includes a formal usability study with domain practitioners, expansion to additional regulatory jurisdictions, and exploration of domain-specific LLM fine-tuning.
The complete source code is available at \url{https://github.com/yuyongkim/chemIP} under the AGPL-3.0 license.

% ============================================================
\section*{Declaration of Competing Interest}

The author declares no competing interests.

\section*{Acknowledgments}

This work utilized public APIs provided by the Korea Occupational Safety and Health Agency (KOSHA), Korea Intellectual Property Rights Information Service (KIPRIS), Korea Trade-Investment Promotion Agency (KOTRA), Ministry of Food and Drug Safety (MFDS), National Institute of Chemical Safety (KISCHEM), and Korea Environment Corporation (NCIS) of the Republic of Korea.
The author thanks the respective agencies for maintaining open data services that enable third-party integration.

% ============================================================
\begin{thebibliography}{99}

\bibitem{occupational_safety_2023}
Korea Occupational Safety and Health Agency (KOSHA),
MSDS OpenAPI Service Technical Documentation, 2024.
\url{https://msds.kosha.or.kr}

\bibitem{pubchem_2023}
S.~Kim, J.~Chen, T.~Cheng, et~al.,
PubChem 2023 update,
Nucleic Acids Research 51~(D1) (2023) D1373--D1380.

\bibitem{scifinder}
Chemical Abstracts Service,
SciFinder-n: Chemical Research Platform, 2024.
\url{https://scifinder-n.cas.org}

\bibitem{reaxys}
Elsevier,
Reaxys: Chemistry Research Platform, 2024.
\url{https://www.reaxys.com}

\bibitem{openfda}
U.S.\ Food and Drug Administration,
openFDA API Documentation, 2024.
\url{https://open.fda.gov}

\bibitem{ahocorasick}
A.V.~Aho, M.J.~Corasick,
Efficient string matching: An aid to bibliographic search,
Communications of the ACM 18~(6) (1975) 333--340.

\bibitem{sqlite_fts5}
D.R.~Hipp,
SQLite FTS5 Extension Documentation, 2024.
\url{https://www.sqlite.org/fts5.html}

\bibitem{ollama}
Ollama,
Run Large Language Models Locally, 2024.
\url{https://ollama.com}

\bibitem{ghs_rev10}
United Nations Economic Commission for Europe,
Globally Harmonized System of Classification and Labelling of Chemicals (GHS Rev.~10), 2023.

\bibitem{fastapi}
S.~Ram\'{i}rez,
FastAPI: Modern, Fast Web Framework for Building APIs with Python,
\url{https://fastapi.tiangolo.com}, 2024.

\bibitem{nextjs}
Vercel,
Next.js: The React Framework for the Web,
\url{https://nextjs.org}, 2025.

\end{thebibliography}

\end{document}
